\chapter[Lezione 16 --- 29 ottobre 2026]{Lezione 16: limiti per sostituzione}
\label{chap:lezione-16}

\section{Teorema sui limiti per sostituzione}
\label{sec:lez16-limiti-sostituzione}

\begin{theorem}[Limite per sostituzione]
\label{thm:lez16-limite-sostituzione}
Sia \(f\colon A\longrightarrow\R\), con
\(A\subseteq\R\) non vuoto, e sia
\(x_0\in\overline{\R}\) un punto di accumulazione per \(A\).
Si supponga che
\[
  \lim_{x\to x_0}f(x)=\ell\in\overline{\R}.
\]
Sia inoltre \(\varphi\colon B\longrightarrow A\), con
\(B\subseteq\R\) non vuoto, e sia
\(t_0\in\overline{\R}\) un punto di accumulazione per \(B\),
tale che
\[
  \lim_{t\to t_0}\varphi(t)=x_0.
\]
Se vale almeno una delle seguenti condizioni:
\begin{enumerate}
  \item \(\varphi(t)\ne x_0\) per ogni
  \(t\in B\setminus\{t_0\}\);
  \item \(x_0\in A\) e \(f(x_0)=\ell\in\R\);
\end{enumerate}
allora
\[
  \lim_{t\to t_0}f(\varphi(t))=\ell.
\]
\end{theorem}

\begin{proof}
Si supponga anzitutto che \(\ell\in\R\).
Sia \(\varepsilon>0\). Per la \cref{def:lez08-limite},
esiste un intorno \(U_{x_0}\) di \(x_0\) tale che
\[
  |f(x)-\ell|<\varepsilon
  \qquad\forall x\in(U_{x_0}\cap A)\setminus\{x_0\}.
\]
Poiché \(\varphi(t)\to x_0\) per \(t\to t_0\),
per la stessa definizione esiste un intorno \(W_{t_0}\)
di \(t_0\) tale che
\[
  \varphi(t)\in U_{x_0}
  \qquad\forall t\in(W_{t_0}\cap B)\setminus\{t_0\}.
\]
Essendo \(\varphi(B)\subseteq A\), ne segue che
\[
  \varphi(t)\in U_{x_0}\cap A
  \qquad\forall t\in(W_{t_0}\cap B)\setminus\{t_0\}.
\]

\emph{Prima condizione.}
Se \(\varphi(t)\ne x_0\) per ogni
\(t\in B\setminus\{t_0\}\), allora
\[
  \varphi(t)\in(U_{x_0}\cap A)\setminus\{x_0\}
  \qquad\forall t\in(W_{t_0}\cap B)\setminus\{t_0\}.
\]
Pertanto
\[
  |f(\varphi(t))-\ell|<\varepsilon
  \qquad\forall t\in(W_{t_0}\cap B)\setminus\{t_0\}.
\]
Per la \cref{def:lez08-limite}, si conclude
\[
  \lim_{t\to t_0}f(\varphi(t))=\ell.
\]

\emph{Seconda condizione.}
Se \(x_0\in A\) e \(f(x_0)=\ell\), la disuguaglianza
precedente vale anche per \(x=x_0\). Pertanto
\[
  |f(x)-\ell|<\varepsilon
  \qquad\forall x\in U_{x_0}\cap A.
\]
Poiché \(\varphi(t)\in U_{x_0}\cap A\) per ogni
\(t\in(W_{t_0}\cap B)\setminus\{t_0\}\), si ha
\[
  |f(\varphi(t))-\ell|<\varepsilon
  \qquad\forall t\in(W_{t_0}\cap B)\setminus\{t_0\},
\]
e quindi
\[
  \lim_{t\to t_0}f(\varphi(t))=\ell.
\]

Si consideri ora \(\ell=+\infty\). In questo caso
la seconda condizione è esclusa, poiché richiede
\(\ell\in\R\); deve quindi valere la prima.
Sia \(K\in\R\). Per la \cref{def:lez08-limite},
esiste un intorno \(U_{x_0}\) tale che
\[
  f(x)>K
  \qquad\forall x\in(U_{x_0}\cap A)\setminus\{x_0\}.
\]
Dal limite di \(\varphi\) si ricava un intorno
\(W_{t_0}\) tale che
\[
  \varphi(t)\in U_{x_0}\cap A
  \qquad\forall t\in(W_{t_0}\cap B)\setminus\{t_0\}.
\]
La prima condizione assicura inoltre che
\(\varphi(t)\ne x_0\) per tali valori di \(t\).
Pertanto
\[
  f(\varphi(t))>K
  \qquad\forall t\in(W_{t_0}\cap B)\setminus\{t_0\},
\]
da cui
\[
  \lim_{t\to t_0}f(\varphi(t))=+\infty.
\]

Infine, se \(\ell=-\infty\), vale necessariamente
la prima condizione. Fissato \(K\in\R\), esiste
un intorno \(U_{x_0}\) tale che
\[
  f(x)<K
  \qquad\forall x\in(U_{x_0}\cap A)\setminus\{x_0\}.
\]
Per il limite di \(\varphi\), esiste un intorno
\(W_{t_0}\) tale che
\[
  \varphi(t)\in U_{x_0}\cap A
  \qquad\forall t\in(W_{t_0}\cap B)\setminus\{t_0\}.
\]
Poiché \(\varphi(t)\ne x_0\), si ottiene
\[
  f(\varphi(t))<K
  \qquad\forall t\in(W_{t_0}\cap B)\setminus\{t_0\},
\]
e quindi
\[
  \lim_{t\to t_0}f(\varphi(t))=-\infty.
\]
\end{proof}

\begin{example}[Importanza delle ipotesi del teorema]
\label{ex:lez16-sostituzione-ipotesi}
Si considerino le funzioni
\[
  f\colon\R\longrightarrow\R,
  \qquad
  f(x):=
  \begin{cases}
    0, & x\ne0,\\
    1, & x=0,
  \end{cases}
\]
e
\[
  \varphi\colon\R\longrightarrow\R,
  \qquad \varphi(t):=0.
\]
Ponendo \(x_0=t_0=0\), si ha
\[
  \lim_{x\to0}f(x)=0,
  \qquad
  \lim_{t\to0}\varphi(t)=0.
\]
Tuttavia, \(\varphi(t)=x_0\) per ogni \(t\in\R\),
perciò la prima condizione del
\cref{thm:lez16-limite-sostituzione} non è soddisfatta.
Neppure la seconda è soddisfatta, poiché
\(f(0)=1\ne0\). Infatti,
\[
  f(\varphi(t))=f(0)=1
  \qquad\forall t\in\R,
\]
da cui
\[
  \lim_{t\to0}f(\varphi(t))=1
  \ne0=\lim_{x\to0}f(x).
\]
\end{example}

\begin{example}[Limite di un'esponenziale per sostituzione]
\label{ex:lez16-sostituzione-esponenziale}
Si vuole calcolare
\[
  \lim_{t\to0}2^{1/t^2}.
\]
Si considerino le funzioni
\[
  f\colon\R\longrightarrow\R,
  \qquad f(x):=2^x,
\]
e
\[
  \varphi\colon\R\setminus\{0\}\longrightarrow\R,
  \qquad \varphi(t):=\frac1{t^2}.
\]
Per il \cref{thm:lez14-limite-rapporto}, poiché
\(t^2\to0\) e \(t^2>0\) per \(t\ne0\), si ha
\[
  \lim_{t\to0}\varphi(t)
  =\lim_{t\to0}\frac1{t^2}=+\infty.
\]
Il punto a cui tende \(\varphi(t)\) è dunque
\(x_0=+\infty\). Per il
\cref{thm:lez15-limiti-esponenziale-logaritmo},
\[
  \lim_{x\to+\infty}f(x)
  =\lim_{x\to+\infty}2^x=+\infty.
\]
La seconda condizione del
\cref{thm:lez16-limite-sostituzione} non può valere,
mentre la prima è soddisfatta poiché
\(\varphi(t)\in\R\) per ogni \(t\ne0\), e quindi
\(\varphi(t)\ne+\infty\).
Per il medesimo teorema, si conclude
\[
  \lim_{t\to0}2^{1/t^2}
  =\lim_{x\to+\infty}2^x=+\infty.
\]
\end{example}
